% Simple poster infographic: Test -> Add controlled error -> Recover -> Measure
% Required packages:
% \usepackage{xcolor}
% \usepackage{graphicx}
% \usepackage{tikz}
% \usetikzlibrary{arrows.meta,positioning}

\definecolor{FlowBlue}{HTML}{1769AA}
\definecolor{FlowPurple}{HTML}{7651B5}
\definecolor{FlowOrange}{HTML}{E56B2F}
\definecolor{FlowGreen}{HTML}{23856D}
\definecolor{FlowInk}{HTML}{172033}

\resizebox{\linewidth}{!}{%
\begin{tikzpicture}[
  font=\sffamily,
  node distance=9mm,
  step/.style={
    rounded corners=4mm,
    draw=#1,
    fill=#1!7,
    line width=1.6pt,
    text width=5.1cm,
    minimum height=6.1cm,
    align=left,
    inner sep=5mm,
    text=FlowInk
  },
  flow/.style={
    -{Stealth[length=4mm,width=3mm]},
    draw=FlowInk!65,
    line width=1.8pt
  }
]

\node[step=FlowBlue] (test) {
  \centering\textcolor{FlowBlue}{\bfseries\Large 1. TEST}\\[3mm]
  \raggedright
  \textbf{8 held-out semiprimes}\\[1.5mm]
  Uniform and discrete-Gaussian Regev samples\\[1.5mm]
  \(M=8,16,32\)\\[1.5mm]
  64 repeats per setting
};

\node[step=FlowPurple,right=of test] (error) {
  \centering\textcolor{FlowPurple}{\bfseries\Large 2. ADD CONTROLLED ERROR}\\[3mm]
  \raggedright
  \textbf{Regev primary test:}\\
  omit 1--2 smallest QFT phase layers\\[1.5mm]
  \textbf{Regev stress tests:}\\
  1\% readout flips, 5\% corrupted shots, and bounded noisy-dual displacement\\[1.5mm]
  \textbf{Shor companion tests:}\\
  phase/RZ rotation error and 0.1--1\% CX depolarizing noise
};

\node[step=FlowOrange,right=of error] (recover) {
  \centering\textcolor{FlowOrange}{\bfseries\Large 3. RECOVER}\\[3mm]
  \raggedright
  Samples\\
  \(\downarrow\)\\
  augmented integer lattice\\
  \(\downarrow\)\\
  LLL reduction\\
  \(\downarrow\)\\
  verify relation + GCD factor
};

\node[step=FlowGreen,right=of recover] (measure) {
  \centering\textcolor{FlowGreen}{\bfseries\Large 4. MEASURE}\\[3mm]
  \raggedright
  Factor-recovery success\\[1.5mm]
  Quantum samples required\\[1.5mm]
  Controlled-phase gates\\[1.5mm]
  CX gates and circuit depth\\[1.5mm]
  LLL runtime
};

\draw[flow] (test) -- (error);
\draw[flow] (error) -- (recover);
\draw[flow] (recover) -- (measure);

\node[
  anchor=north,
  text width=12.8cm,
  align=center,
  font=\sffamily\small,
  text=FlowInk!80
] at ($(error.south east)!0.5!(recover.south west)+(0,-7mm)$) {
  \textbf{Important:} these are controlled simulations.
  QFT truncation is an algorithmic approximation, and the other channels are
  declared noise surrogates---not measurements from real quantum hardware.
};

\end{tikzpicture}%
}
